% Study 1: supplementary methods (materials, procedure, exclusions and sample), the terms its
% tables use, and its tables and figure for the Supplementary Information. Numbers come only from
% macros in generated/numbers.tex and from the Study 1 tables (supplement/tables/s1_*.tex) that
% tools/make_tables.py builds from the pipeline's CSV tables; supplement/tables/README_suppA.md
% lists each table and its source.
\label{sec:suppA-study1}
% Within this section every table and figure is set exactly where it is input ([H], from the float
% package), so that each follows the sentence that introduces it. The tables ask for [htbp], which is
% overridden here without editing them.
\begingroup
\let\suppAtable\table
\renewcommand{\table}[1][]{\suppAtable[H]}

Study~1 asked whether a label saying that an AI system made an image lowers how creative, and how effortful, people judge that same image to be. Its \SOneN{} participants each rated ten images for creativity and for effort. Each participant was told one of two things about the artist, Alex: that he made the images with design and sketching tools (the Human label) or by prompting an AI system (the AI label). The main text reports the registered tests of Study~1's four hypotheses (H1 to H4). Supplementary Note~1 gives H3 and H4 in full, Supplementary Note~2 the robustness checks and the belief check, and Supplementary Note~7 what the written explanations referred to; Supplementary Table~\ref{tab:supp-verdicts} lists each verdict. A registration is a dated record of a study's hypotheses and analysis plan, unrelated to copyright registration. Study~1's was filed on the Open Science Framework (OSF), an online archive for study plans, on 8 October 2026, after its data had been collected and analyzed (Section~\ref{sec:suppA-registrations}). This section gives the materials, the procedure and the sample, defines the terms its tables use, and presents those tables.

\subsection*{Materials}

Every participant saw the same ten images. All ten were generated with OpenAI's DALL·E, an AI image generator, used through ChatGPT and given prompts that we wrote (Section~\ref{sec:suppA-registrations}), so the Human label was a cover story. Five images show a human or humanoid figure and five do not; the analysis calls this the kind of image. The same ten images were the originals of Studies 2 to 5.

Participants read the label once, at the start. The Human label read: `There is a digital artist named Alex. Alex is a skilled artist who creates digital artworks. Alex uses his graphic design skills and digital sketching tools.' The AI label read: `There is a digital artist named Alex. Alex is a skilled artist who creates digital artworks. Alex develops and inputs prompts into an artificial intelligence (AI) system. The AI system generates digital artworks based on those prompts.' The AI label is one sentence longer.

Before each image, the Human group read `Alex had a new idea for a digital artwork. Alex worked with his graphic design and digital sketching tools, revising as his work as he went, which created a digital artwork.' The AI group read `Alex had a new idea for a work of digital artwork. Alex developed and input prompts into the AI image generator, revising each prompt as he went, which created a digital artwork.' Both went on: `Click the arrow to see the digital artwork created. Note that you will be required to stay on the page for at least 8 seconds before proceeding. Please pay close attention to the artwork.' These quotations reproduce the survey exactly, slips included.

Each image was rated twice, on five-point scales scored 1 to 5 by the position of the answer chosen. The creativity item read `How much creativity was involved in the production of this work? Creativity is defined as the presence of at least a minimal original, intellectual contribution to the creation of the work.' Its answers ran from `No creativity involved' through minimal, moderate and significant creativity to `A great deal of creativity involved'. The effort item read `How much effort was involved in the production of this work?' and carried no definition. Its answers ran from `No effort involved' through minimal, moderate and significant effort to `A great deal of effort involved'. One effort item, for the first image with a figure, labeled its top two answers differently: `A lot of effort involved' and `Significant effort involved'. It is scored by position like the others, and a registered check drops it (Supplementary Note~2). After the two ratings of an image, two required text boxes asked `Please explain your reasoning for your rating on creativity:' and `Please explain your reasoning for your rating on effort:'. Copying and pasting into these boxes were not blocked.

\subsection*{Procedure}

By our account, participants were students in the UCLA Psychology subject pool who took part for course credit. Every analyzed record carries the code that the survey sent, at its end, to the pool's page for granting credit. The consent form, the same in all five studies, also offered hourly pay to people who signed up outside the pool, and no record shows that anyone took that route. After the consent form and the instructions, a task check asked `What is your task in this experiment?' with five answers, the correct one being `Viewing images for a set amount of time and answering questions'. An answer was not required. A wrong answer brought the instructions back once, and the participant then went on whatever the second answer.

A randomizer in the survey software, Qualtrics, then assigned each participant to one of the two labels, presenting each equally often. The survey also set a cap of 60 participants per label, but the cap never changed which label anyone saw. The ten images followed in a random order that the survey did not record. For each image, the description above came first and then the image, which stayed on screen for at least 8 seconds. A separate screen, without the image, then asked for the creativity rating and the effort rating, followed by the explanation of each, in that order. Every explanation was therefore written after both ratings of its image. A final block asked `Did you complete this experiment seriously throughout (without responding randomly)?', with the answers `No, I didn't' and `Yes, I did', then whether any part of the procedure had been unclear, and then race, gender, age, native language and formal artistic training. No question asked whether participants noticed or believed the label. The debriefing, the same in all five studies, told participants that `Current legal standards state that anything produced with artificial intelligence is insufficiently creative, even if you use AI as a base against which you create new work.' That sentence overstates the law. The US Copyright Office refuses copyright registration for material whose expressive elements an AI system determined, but it does register a person's own creative selection or arrangement of that material, or a modification of it creative enough to meet copyright's standard of originality, with the AI-generated material left out of the claim.\cite{usco2023guidance,usco2025copyrightability} Participants read the debriefing only after their last rating.

\subsection*{Exclusions and sample}

The survey export held \SOneNRecords{} records. The registered exclusions removed, in order, \SOneNPreviewRemoved{} survey previews (test runs by the researchers), \SOneNUnfinishedRemoved{} unfinished records, and \SOneNNotSeriousRemoved{} records with any answer but `Yes, I did' to the seriousness question, a blank included. The task check removed \SOneNAttentionRemoved{} records, because it let everyone continue. That left \SOneN{} participants, \SOneNHuman{} under the Human label and \SOneNAI{} under the AI label. Each gave all 20 ratings and 20 explanations, \SOneNRatings{} of each in all. The data were collected from 25 February to 3 March 2025. Supplementary Table~\ref{tab:supp_exclusions} gives the flow, and Supplementary Table~\ref{tab:supp_attrition} gives the removals by label, with Fisher's exact test, a test for counts in small tables, of whether the share removed differed between labels.

The participants' mean age was \SOneAgeMean{} years (standard deviation, SD, a measure of spread: \SOneAgeSD; range \SOneAgeMin{} to \SOneAgeMax). Of the \SOneN, \SOneNWomen{} answered the gender question Female, \SOneNMen{} Male and \SOneNGenderOther{} Other. Supplementary Table~\ref{tab:supp_demographics} gives race and formal artistic training as well. Supplementary Table~\ref{tab:supp_balance} compares the two label groups on these background measures, as a check that random assignment produced similar groups.

% Study 1: exclusions.
% Written by tools/make_tables.py from the pipeline table(s) below (paths relative to
% Analysis/R Analysis Pipeline, each with the SHA-256 it was read with); do not edit by hand.
%   Study 1/tables/study1_exclusion_flow.csv  5a76f3fe7f18f46d3c304d1e538b31f3ae1a17975029d55dcb384cf18eab7d10
\begin{table}[htbp]
\centering
\small
\caption[Study 1: exclusions]{\textbf{Study~1: exclusions.} The registered exclusion steps in the order they were applied, with the records each step removed and the records left after it. A record is one session in the survey export.}\label{tab:supp_exclusions}
\setlength{\tabcolsep}{4pt}
\begin{tabular}{@{}>{\raggedright\arraybackslash}p{11.4cm}rr@{}}
\toprule
Step & Removed & Records left\\
\midrule
Records in the survey export, test runs included & -- & 119\\
Survey previews (test runs by the researchers) removed & 2 & 117\\
Unfinished records removed & 3 & 114\\
Any answer but `Yes, I did' to the seriousness question, a blank included, removed & 3 & 111\\
Wrong answer to the task check removed & 0 & 111\\
\bottomrule
\end{tabular}
\par\smallskip
\begin{minipage}{\textwidth}\footnotesize\raggedright The task check asked `What is your task in this experiment?'. A wrong answer brought the instructions back once and then let the participant continue, so this step removed no one. The seriousness question asked `Did you complete this experiment seriously throughout (without responding randomly)?'.\par\end{minipage}
\end{table}

% Study 1: removals by label.
% Written by tools/make_tables.py from the pipeline table(s) below (paths relative to
% Analysis/R Analysis Pipeline, each with the SHA-256 it was read with); do not edit by hand.
%   Study 1/tables/study1_attrition_by_condition.csv  ad43935243eb6c87150c77cd64dab1f7ff34e901a0dbf4cb34a79eb6d1623f6e
\begin{table}[htbp]
\centering
\small
\caption[Study 1: removals by label]{\textbf{Study~1: removals by label.} The records left once survey previews had been removed, by the label each participant had been shown: those analyzed and those removed at each later step (descriptive).}\label{tab:supp_attrition}
\setlength{\tabcolsep}{4pt}
\begin{tabular}{@{}lrrr@{}}
\toprule
Records & Human label & AI label & Total\\
\midrule
Analyzed & 56 & 55 & 111\\
Removed as unfinished & 1 & 2 & 3\\
Removed by the seriousness question & 2 & 1 & 3\\
\bottomrule
\end{tabular}
\par\smallskip
\begin{minipage}{\textwidth}\footnotesize\raggedright Fisher's exact test, a test for counts in small tables, compared the two labels. Removed at any step against analyzed: $p$~\ensuremath{>}~.999. The whole table of steps by label: $p$~\ensuremath{>}~.999. Every record carried a label.\par\end{minipage}
\end{table}

% Study 1: the analyzed participants.
% Written by tools/make_tables.py from the pipeline table(s) below (paths relative to
% Analysis/R Analysis Pipeline, each with the SHA-256 it was read with); do not edit by hand.
%   Study 1/tables/study1_demo_age.csv  95fda01292f4a4cc41845ddc3de73c77a497c3e65747bb758ba236d7a11118a4
%   Study 1/tables/study1_demo_gender.csv  846ba92df512e0f199472d00abc578f58ff9ab5470084edd7a93238b3d3465e5
%   Study 1/tables/study1_demo_race.csv  dac3c3e28a8942ff391aab00e9c04a2dcaf1cfe3b6e64907d9c2f22e61329b5d
%   Study 1/tables/study1_demo_experience.csv  590883417fee4e75547b3369e3026420a22a6e2e4b465996c49b0aeab2d0bf03
%   Study 1/tables/study1_demo_balance.csv  1c1a883be1c80d7a1de4364e5a61b2e8946a271a2b2973cc4cc33bb92cd5cbfb
\begin{table}[htbp]
\centering
\small
\caption[Study 1: the analyzed participants]{\textbf{Study~1: the analyzed participants.} Their ages, their answers to the background questions at the end of the survey, and the label each was shown (descriptive).}\label{tab:supp_demographics}
\setlength{\tabcolsep}{4pt}
\begin{tabular}{@{}>{\raggedright\arraybackslash}p{3.6cm}>{\raggedright\arraybackslash}p{8.9cm}r@{}}
\toprule
Question & Answer & Participants (\%)\\
\midrule
Age in years & Mean 20.08 (SD 1.69); range 18 to 29 & 111 (100.0\%)\\
\addlinespace[3pt]
Gender & Female & 93 (83.8\%)\\
 & Male & 17 (15.3\%)\\
 & Other & 1 (0.9\%)\\
\addlinespace[3pt]
Race & Asian & 35 (31.5\%)\\
 & White or Caucasian & 35 (31.5\%)\\
 & Other & 22 (19.8\%)\\
 & White or Caucasian; Asian & 6 (5.4\%)\\
 & Black or African American & 4 (3.6\%)\\
 & White or Caucasian; Black or African American & 3 (2.7\%)\\
 & American Indian/Native American or Alaska Native & 1 (0.9\%)\\
 & Asian; Native Hawaiian or Other Pacific Islander & 1 (0.9\%)\\
 & Black or African American; Other & 1 (0.9\%)\\
 & Prefer not to say & 1 (0.9\%)\\
 & White or Caucasian; Asian; Native Hawaiian or Other Pacific Islander & 1 (0.9\%)\\
 & White or Caucasian; Other & 1 (0.9\%)\\
\addlinespace[3pt]
Formal artistic training & No formal training or experience & 43 (38.7\%)\\
 & A little formal training or experience & 41 (36.9\%)\\
 & Moderate formal training or experience & 21 (18.9\%)\\
 & Extensive formal training or experience & 5 (4.5\%)\\
 & Very extensive formal training or experience & 1 (0.9\%)\\
\addlinespace[3pt]
Label shown & Human label & 56 (50.5\%)\\
 & AI label & 55 (49.5\%)\\
\bottomrule
\end{tabular}
\par\smallskip
\begin{minipage}{\textwidth}\footnotesize\raggedright SD, standard deviation, a measure of spread. Percentages are of all analyzed participants. Answers are in the survey's own words. The gender question offered Male, Female, Other and Prefer not to say. Participants could choose more than one race; a row that names several answers, separated by semicolons, counts those who chose exactly that combination. Race was self-reported, by choosing one or more of the survey's categories; it is reported to describe the sample and to check that the label groups were alike (Supplementary Table~\ref{tab:supp_balance}), and no analysis compared racial groups. The training question asked for the level of `formal artistic training or experience, such as painting, drawing, sculpture, ceramics, photography, and/or film-video'.\par\end{minipage}
\end{table}

% Study 1: background of the two label groups.
% Written by tools/make_tables.py from the pipeline table(s) below (paths relative to
% Analysis/R Analysis Pipeline, each with the SHA-256 it was read with); do not edit by hand.
%   Study 1/tables/study1_balance.csv  794c8844a3e4f130cbd245beb28aa7ed1c920da9ba1496080475832220f05ca9
\begin{table}[htbp]
\centering
\small
\caption[Study 1: background of the two label groups]{\textbf{Study~1: did the two label groups differ in background?} Tests comparing the participants shown the Human label with those shown the AI label, as a check that random assignment produced similar groups (descriptive; the four tests are not corrected).}\label{tab:supp_balance}
\setlength{\tabcolsep}{4pt}
\begin{tabular}{@{}>{\raggedright\arraybackslash}p{7.0cm}>{\raggedright\arraybackslash}p{4.0cm}ll@{}}
\toprule
Background measure & Test & Statistic & $p$\\
\midrule
Age & Welch's $t$ test & $t$(92.1)~=~1.88 & $p$~\ensuremath{=}~.064\\
Gender (female, male, other) & Fisher's exact test & -- & $p$~\ensuremath{=}~.519\\
Formal artistic training (five ordered answers) & Wilcoxon rank-sum test & $W$~=~1{,}579 & $p$~\ensuremath{=}~.809\\
Race (the four most common answers; all others pooled) & Fisher's exact test & -- & $p$~\ensuremath{=}~.992\\
\bottomrule
\end{tabular}
\par\smallskip
\begin{minipage}{\textwidth}\footnotesize\raggedright Welch's $t$ test compares two means without assuming that the two groups vary equally; its degrees of freedom are in parentheses. Fisher's exact test is a test for counts in small tables; for race, its $p$-value was simulated from 10{,}000 random tables. The Wilcoxon rank-sum test compares two groups' ordered answers by their ranks; $W$ is its statistic. Every analyzed participant enters every test.\par\end{minipage}
\end{table}


\subsection*{Terms used in the tables}

The tables below use the paper's own terms, and each table's note defines its abbreviations. A few terms recur. Most tests rest on a linear mixed model, a regression that predicts each rating from the label, the kind of image and the two together, while giving each participant and each image a baseline of its own. A random slope lets an effect also vary from one participant, or one image, to the next; a `maximal' structure includes every random slope the design allows. A specification curve refits the same quantity under every combination of a set of reasonable analysis choices. An ordinal model treats the five answers as ordered categories rather than as evenly spaced numbers, and it reports its estimates in log-odds, the logarithm of the odds of a higher answer.

The H1 contrast is the label's penalty: the Human-label mean minus the AI-label mean, averaged with equal weight over the two kinds of image. The primary model's AI-label coefficient is a different quantity. It is the AI-minus-Human difference for images without a figure only, so its sign runs the other way.

A $p$-value is the chance of a result at least this far from no effect if there were no effect. A 95\% confidence interval (CI) is the range of values the data do not rule out at the 95\% level. A fit converged when it settled on a solution, and it is singular when it estimated a variance or a correlation at the edge of its possible range. File names ending in .csv, such as \nolinkurl{study1_specification_curve.csv}, name the full tables that the analysis code writes to Study~1's tables folder.

\subsection*{Results not reported elsewhere}

\paragraph*{The confirmatory tests and the models behind them.}
Study~1's registration named four confirmatory tests and corrected them together by Holm's method,\cite{holm1979simple} which raises each $p$-value so that the chance of calling any of the four significant by mistake stays at .05. Supplementary Table~\ref{tab:supp_holm_family} lists the four with their $p$-values before and after the correction. Its H4 row is the test of spread between people on creativity alone, and passing it does not mean that H4 passed (Supplementary Note~1). Supplementary Table~\ref{tab:supp_lmm} gives every estimate of the primary model of each rating. Supplementary Table~\ref{tab:suppA-h2} gives H2 in its two registered versions, which agree.

% Study 1: the four confirmatory tests.
% Written by tools/make_tables.py from the pipeline table(s) below (paths relative to
% Analysis/R Analysis Pipeline, each with the SHA-256 it was read with); do not edit by hand.
%   Study 1/tables/study1_confirmatory_family.csv  d502d85829e95b88558fc42e57ea4256f2c2506ef932396cfb090241b2d06fbc
%   Study 1/tables/study1_joint_model_anova.csv  c7e91d3bef09735c81a4b0b0533c03e9c6db850983f615b65775cfbea7543c83
%   Study 1/tables/study1_h4_variance_test.csv  499b8678ab6a36e9b9d5b1d9bfa2c9b547b16020d745fde8a443cecb28b62e07
\begin{table}[htbp]
\centering
\small
\caption[Study 1: the four confirmatory tests]{\textbf{Study~1: the four confirmatory tests and Holm's correction.} The registration named these four as its confirmatory tests and corrected them together by Holm's method (registered confirmatory). All four pass after the correction. H1 and H2 are supported; H4 is not, because it needed more than this one test (note a).}\label{tab:supp_holm_family}
\setlength{\tabcolsep}{4pt}
\begin{tabular}{@{}>{\raggedright\arraybackslash}p{6.4cm}llll@{}}
\toprule
Test & Estimate (SE) & Statistic & $p$ & Holm $p$\\
\midrule
H1, creativity: the label's penalty & 0.45 (0.13) & $t$(109)~=~3.49 & $p$~\ensuremath{<}~.001 & $p$~\ensuremath{=}~.001\\
H1, effort: the label's penalty & 0.92 (0.15) & $t$(109)~=~6.08 & $p$~\ensuremath{<}~.001 & $p$~\ensuremath{<}~.001\\
H2: the penalty on effort minus the penalty on creativity & 0.47 (0.10) & $F$(1,~109)~=~23.02 & $p$~\ensuremath{<}~.001 & $p$~\ensuremath{<}~.001\\
H4, between-person part, creativity: AI-to-Human ratio of the variance between people\textsuperscript{a} & 2.04 & $\chi^2$(1)~=~4.98 & $p$~\ensuremath{=}~.026 & $p$~\ensuremath{=}~.026\\
\bottomrule
\end{tabular}
\par\smallskip
\begin{minipage}{\textwidth}\footnotesize\raggedright The penalty is the Human-label mean rating minus the AI-label mean rating, averaged with equal weight over images with and without a human figure, in points on the 1-to-5 scale. SE, standard error. Degrees of freedom are in parentheses. Holm's method orders the four $p$-values from smallest to largest and compares each with a threshold that starts strict and relaxes as tests pass, so that the chance of calling any of the four significant by mistake stays at .05; the Holm $p$ is each $p$-value adjusted in that way. The H4 row is a likelihood-ratio test, which compares a model in which the variance between people differs by label with one in which it does not. It is shown for creativity, the weaker of the two ratings, and its estimate has no standard error. \textsuperscript{a}This row is only the between-person part of H4. H4 as registered also required a smaller variance within people under the AI label, with both ratios' resampled intervals on the predicted side of 1, on both ratings. That held on effort and not on creativity, so H4 was not supported (Supplementary Note~1).\par\end{minipage}
\end{table}

% Study 1: every estimate of the primary model.
% Written by tools/make_tables.py from the pipeline table(s) below (paths relative to
% Analysis/R Analysis Pipeline, each with the SHA-256 it was read with); do not edit by hand.
%   Study 1/tables/study1_lmm_fixed_effects.csv  e050698a9459c7df8aede2cdd6c34fd54ce2c5d2203b53be11a3ec5e25eeb942
\begin{table}[htbp]
\centering
\small
\caption[Study 1: every estimate of the primary model]{\textbf{Study~1: every estimate of the primary model of each rating.} The registered model, fitted to each rating separately; its coefficients are descriptive. The registered H1 test is not the AI-label row here, which is the effect for images without a figure only and runs AI minus Human, but the penalty averaged over both kinds of image (Supplementary Table~\ref{tab:supp_holm_family}).}\label{tab:supp_lmm}
\setlength{\tabcolsep}{4pt}
\begin{tabular}{@{}>{\raggedright\arraybackslash}p{8.8cm}lll@{}}
\toprule
Term & Estimate (SE) & 95\% CI & $p$\\
\midrule
\multicolumn{4}{@{}l}{\textit{Creativity}}\\
Baseline: Human label, images without a figure & 3.37 (0.22) & 2.88 to 3.86 & $p$~\ensuremath{<}~.001\\
AI label, images without a figure (AI minus Human) & \ensuremath{-}0.41 (0.14) & \ensuremath{-}0.68 to \ensuremath{-}0.14 & $p$~\ensuremath{=}~.003\\
Images with a figure, Human label (with minus without) & 0.35 (0.29) & \ensuremath{-}0.32 to 1.02 & $p$~\ensuremath{=}~.268\\
Label by kind of image (the H3 quantity) & \ensuremath{-}0.08 (0.09) & \ensuremath{-}0.27 to 0.11 & $p$~\ensuremath{=}~.393\\
\addlinespace[3pt]
\multicolumn{4}{@{}l}{\textit{Effort}}\\
Baseline: Human label, images without a figure & 3.72 (0.18) & 3.34 to 4.10 & $p$~\ensuremath{<}~.001\\
AI label, images without a figure (AI minus Human) & \ensuremath{-}0.91 (0.16) & \ensuremath{-}1.22 to \ensuremath{-}0.60 & $p$~\ensuremath{<}~.001\\
Images with a figure, Human label (with minus without) & 0.09 (0.21) & \ensuremath{-}0.39 to 0.57 & $p$~\ensuremath{=}~.685\\
Label by kind of image (the H3 quantity) & \ensuremath{-}0.02 (0.10) & \ensuremath{-}0.21 to 0.17 & $p$~\ensuremath{=}~.833\\
\bottomrule
\end{tabular}
\par\smallskip
\begin{minipage}{\textwidth}\footnotesize\raggedright Estimates are in points on the 1-to-5 scale. The model is a linear mixed model, which gives each participant and each image a baseline of its own; the Human label and images without a human figure are its reference levels. The label-by-kind-of-image row is how much the AI label's effect changes when the image has a figure, the quantity H3 tests. SE, standard error; CI, confidence interval. The $p$-values use Satterthwaite's approximation to the degrees of freedom.\par\end{minipage}
\end{table}

% Study 1's H2 in its two registered versions, every number a macro from generated/numbers.tex
% (sources: Study 1/tables/study1_joint_model_contrast_of_contrasts.csv, study1_confirmatory_family.csv
% and study1_h2_ordinal.csv). See supplement/tables/README_suppA.md.
\begin{table}[tbp]
\centering
\small
\caption{\textbf{Study~1's H2 in its two registered versions.} H2 predicted that the AI label lowers effort ratings more than creativity ratings. The quantity is the label's penalty on effort minus its penalty on creativity, averaged over the two kinds of image, where a penalty is the Human-label mean minus the AI-label mean. Both versions fit both ratings together, with a baseline and an effort-minus-creativity difference for each participant and a baseline for each image on each rating. The linear model treats the five answers as evenly spaced numbers, so its estimate is in rating points. The ordinal model treats them as ordered categories, so its estimate is on that model's log-odds scale and cannot be compared in size with the first. The registration requires the two to agree in sign and significance; they do, so H2 is supported (registered confirmatory). Neither fit was singular. CI, confidence interval.}
\label{tab:suppA-h2}
\begin{tabular}{@{}>{\raggedright\arraybackslash}p{3.6cm}lllll@{}}
\toprule
Version & Estimate & 95\% CI & SE & Test statistic & $p$\\
\midrule
Linear mixed model (rating points) & \SOneHTwoDiff & \SOneHTwoLo{} to \SOneHTwoHi & \SOneHTwoSE & $t$(\SOneHTwoDf)~=~\SOneHTwoT & $p$~\SOneHTwoP; Holm $p$~\SOneHTwoPHolm\\
Ordinal mixed model (log-odds) & \SOneHTwoOrdEst & -- & \SOneHTwoOrdSE & $z$~=~\SOneHTwoOrdZ & $p$~\SOneHTwoOrdP\\
\bottomrule
\end{tabular}
\end{table}


\paragraph*{H3 and H4.}
Supplementary Note~1 reports both hypotheses, with their estimates, intervals, power and verdicts. It does not print H4's likelihood-ratio tests. Each compares a model in which one spread differs between the label groups with a model in which it does not, and its statistic is referred to a chi-square ($\chi^2$) distribution with one degree of freedom. For the spread between people, the tests gave $\chi^2$(1)~=~\SOneHFourCreatChisq{} on creativity ($p$~\SOneHFourCreatP) and $\chi^2$(1)~=~\SOneHFourEffChisq{} on effort ($p$~\SOneHFourEffP); for the spread within people, $p$~\SOneHFourCreatResidualP{} and $p$~\SOneHFourEffResidualP. H4's verdict does not rest on these tests. It rests on the two ratios of spreads and their resampled intervals, and as registered H4 was not supported (Supplementary Note~1). In simulated samples of this size, the estimated ratio of spreads between people ranged across its middle 95\% from \SOneHFourCreatResolutionRatioLo{} to \SOneHFourCreatResolutionRatioHi{} on creativity and from \SOneHFourEffResolutionRatioLo{} to \SOneHFourEffResolutionRatioHi{} on effort, when the true ratios equal those observed (registered description of what the design can resolve). The analysis tables \nolinkurl{study1_h3_equivalence.csv}, \nolinkurl{study1_h3_equivalence_power.csv}, \nolinkurl{study1_h4_variance_test.csv} and \nolinkurl{study1_h4_resolution.csv} give every row.

\paragraph*{Robustness of the penalty.}
Supplementary Note~2 describes these registered checks and their results. Supplementary Table~\ref{tab:supp_table_1} gathers several of them; its specification-curve row summarizes the \SOneSpecNTotal{} refits of the specification curve, each of which is in the analysis table \nolinkurl{study1_specification_curve.csv}. The table's Bayesian rows report the primary model's AI-label coefficient, not the H1 contrast. The Bayesian refits estimate the model by combining the data with prior assumptions and summarize the result as a distribution. In them, the H1 contrast was \SOneBayesCreat{} points on creativity (95\% credible interval, the range holding 95\% of that distribution, \SOneBayesCreatLo{} to \SOneBayesCreatHi) and \SOneBayesEff{} on effort (\SOneBayesEffLo{} to \SOneBayesEffHi). Supplementary Fig.~\ref{fig:suppA-ppc} checks the linear Bayesian refit, which treats ratings as continuous numbers as the primary model does, against the data. The registered maximal structure lets the label's penalty vary from image to image and the effect of the kind of image vary from person to person. With it, the H1 contrast was \SOneMaximalCreatEst{} on creativity ($p$~\SOneMaximalCreatP) and \SOneMaximalEffEst{} on effort ($p$~\SOneMaximalEffP). Both fits were singular, so these values are reported with that status (analysis table \nolinkurl{study1_lmm_fully_maximal.csv}).

\begin{figure}[H]
\centering
% Drawn by tools/make_figures.py from figures/figS_study1_ppc_curves.csv, which
% tools/study1_ppc_data.R computes from the pipeline's saved Bayesian fits.
\makebox[\textwidth][c]{\includegraphics[width=183mm]{figS_study1_ppc}}
\caption{\textbf{Study~1: the linear Bayesian refit checked against the data.} A posterior predictive check compares ratings simulated from a fitted model with the ratings observed: \textbf{a},~creativity; \textbf{b},~effort. The horizontal axis is the rating; the height of a curve shows how common each rating value is, smoothed across neighboring values. The black line is the observed ratings; the gray lines are 100 sets of ratings simulated from the saved model, with the draws fixed by a seed. The model treats ratings as continuous numbers, so its simulated ratings are smooth and run past 1 and 5, while the observed ratings take only the five scale values. This mismatch is why the registration also asks for the ordinal model, which treats the answers as ordered categories. The Bayesian refit is a registered secondary check.}
\label{fig:suppA-ppc}
\end{figure}

% Study 1: robustness checks on the penalty.
% Written by tools/make_tables.py from the pipeline table(s) below (paths relative to
% Analysis/R Analysis Pipeline, each with the SHA-256 it was read with); do not edit by hand.
%   Study 1/tables/study1_confirmatory_family.csv  d502d85829e95b88558fc42e57ea4256f2c2506ef932396cfb090241b2d06fbc
%   Study 1/tables/study1_specification_curve.csv  bf372ede20c3a9887df35cbeacdeda5983fe44a4d0c7f9e7dc06a8325ecf4c09
%   Study 1/tables/study1_brms_primary_fixed.csv  60514547f37adf2c52b47d6e914a82ef2db0380e86e9f7caf49028948101ef5b
%   Study 1/tables/study1_brms_cumulative_fixed.csv  865874e9bead54ac27d9f389a0b49f6c01badcfd61c5c84b962a465b267cc34c
%   Study 1/tables/study1_lmm_fixed_effects.csv  e050698a9459c7df8aede2cdd6c34fd54ce2c5d2203b53be11a3ec5e25eeb942
%   Study 1/tables/study1_equivalence_sensitivity.csv  be48d50f8987f8042a6e995e697c338c383eb78fa584ea70b1ae1727dbd1f634
\begin{table}[htbp]
\centering
\footnotesize
\caption[Study 1: robustness checks on the penalty]{\textbf{Study~1: robustness checks on the label's penalty.} Each row refits or re-tests the H1 contrast, or the related quantity the row names, under another analysis that the registration lists. The status of each row is in parentheses; none of these checks carries a verdict of its own.}\label{tab:supp_table_1}
\setlength{\tabcolsep}{4pt}
\begin{tabular}{@{}>{\raggedright\arraybackslash}p{4.4cm}>{\raggedright\arraybackslash}p{5.15cm}>{\raggedright\arraybackslash}p{5.15cm}@{}}
\toprule
Check (status) & Creativity & Effort\\
\midrule
The H1 contrast, primary model \newline \textit{(registered confirmatory)} & 0.45 (SE 0.13); $t$(109)~=~3.49, $p$~\ensuremath{<}~.001; Holm $p$~\ensuremath{=}~.001 & 0.92 (SE 0.15); $t$(109)~=~6.08, $p$~\ensuremath{<}~.001; Holm $p$~\ensuremath{<}~.001\\
\addlinespace[4pt]
Specification curve: the H1 contrast refitted under every combination of the registered analysis choices \newline \textit{(registered secondary; the registration calls it descriptive)} & 18 of 18 specifications estimated, 18 recorded as converged; 18 positive and 0 negative; 18 with $p$ below .05; median $p$~\ensuremath{<}~.001 & 18 of 18 specifications estimated, 18 recorded as converged; 18 positive and 0 negative; 18 with $p$ below .05; median $p$~\ensuremath{<}~.001\\
\addlinespace[4pt]
Bayesian refit, linear: the AI-label effect for images without a figure (AI minus Human) \newline \textit{(registered secondary)} & Posterior mean \ensuremath{-}0.40 (95\% credible interval \ensuremath{-}0.66 to \ensuremath{-}0.13); 0.01 points from the primary model's estimate & Posterior mean \ensuremath{-}0.90 (95\% credible interval \ensuremath{-}1.21 to \ensuremath{-}0.59); 0.01 points from the primary model's estimate\\
\addlinespace[4pt]
Bayesian refit, ordinal: the same effect, in log-odds \newline \textit{(registered secondary)} & Posterior mean \ensuremath{-}0.93 (95\% credible interval \ensuremath{-}1.57 to \ensuremath{-}0.29) & Posterior mean \ensuremath{-}1.92 (95\% credible interval \ensuremath{-}2.66 to \ensuremath{-}1.15)\\
\addlinespace[4pt]
Equivalence at wider bounds: the label-by-kind-of-image interaction, the H3 quantity \newline \textit{(descriptive; no verdict on H3)} & \ensuremath{-}0.08 (90\% CI \ensuremath{-}0.24 to 0.08); inside $\pm$0.30 ($p$~\ensuremath{=}~.010) and inside $\pm$0.50 ($p$~\ensuremath{<}~.001) & \ensuremath{-}0.02 (90\% CI \ensuremath{-}0.18 to 0.14); inside $\pm$0.30 ($p$~\ensuremath{=}~.002) and inside $\pm$0.50 ($p$~\ensuremath{<}~.001)\\
\bottomrule
\end{tabular}
\par\smallskip
\begin{minipage}{\textwidth}\footnotesize\raggedright The H1 contrast is the Human-label mean rating minus the AI-label mean rating, averaged with equal weight over images with and without a human figure, in points on the 1-to-5 scale. SE, standard error; CI, confidence interval; $\pm$, plus or minus, in rating points. The specification curve crosses the registered choices of random effects, model family, covariates and level of aggregation; its estimates are on each family's own scale, so only their signs are counted. A Bayesian refit estimates the model by combining the data with prior assumptions and summarizes each quantity as a distribution; the credible interval holds the middle 95\% of it. The two Bayesian rows report the AI-label coefficient of Supplementary Table~\ref{tab:supp_lmm}, not the H1 contrast; the text gives the H1 contrast from the same refits. An ordinal model treats the five answers as ordered categories, and log-odds are the logarithm of the odds of a higher answer. An equivalence test asks whether an effect is shown to lie inside a bound; its $p$ is the larger of two one-sided tests. The last row applies such tests to the H3 quantity at wider bounds than the registered one, as description, and carries no verdict on H3; the registered test of H3, at $\pm$0.20, is in Supplementary Note~1.\par\end{minipage}
\end{table}


Supplementary Table~\ref{tab:supp_influence} gives the largest and the average change in the penalty when the model is refitted once without each participant and once without each image. Supplementary Table~\ref{tab:supp_item_sensitivity} gives H1 on effort and H2 with the differently labeled effort item set aside, and H2 with that image set aside from both ratings. The registration names this check as a sensitivity analysis, settled after the estimates had been read (registered secondary).

% Study 1: leaving out one participant or one image.
% Written by tools/make_tables.py from the pipeline table(s) below (paths relative to
% Analysis/R Analysis Pipeline, each with the SHA-256 it was read with); do not edit by hand.
%   Study 1/tables/study1_influence_lopo.csv  199b870aa85ac5a1074944aaac070a5fcd4287ea7505a961b225e6ad83ab6e75
%   Study 1/tables/study1_influence_loso.csv  c17b0224327a1532aff03995fa274b44509a1d136bb408aef06f4f3d23341e13
\begin{table}[htbp]
\centering
\small
\caption[Study 1: leaving out one participant or one image]{\textbf{Study~1: does any one participant or image carry the penalty?} The primary model was refitted once without each participant and once without each image (registered secondary). The table gives the largest and the average size of the resulting change in the H1 contrast. Every refit kept the penalty above zero.}\label{tab:supp_influence}
\setlength{\tabcolsep}{4pt}
\begin{tabular}{@{}llrrrr@{}}
\toprule
Left out in turn & Rating & Refits & Usable & Largest change & Average change\\
\midrule
Each participant & Creativity & 111 & 111 & 0.03 & 0.01\\
 & Effort & 111 & 111 & 0.04 & 0.01\\
\addlinespace[3pt]
Each image & Creativity & 10 & 10 & 0.04 & 0.01\\
 & Effort & 10 & 10 & 0.03 & 0.01\\
\bottomrule
\end{tabular}
\par\smallskip
\begin{minipage}{\textwidth}\footnotesize\raggedright The H1 contrast is the Human-label mean rating minus the AI-label mean rating, averaged with equal weight over images with and without a human figure. Changes are in points on the 1-to-5 scale, without regard to sign. A usable refit is one that converged without error. The analysis tables \nolinkurl{study1_influence_lopo.csv} and \nolinkurl{study1_influence_loso.csv} give the estimate from every refit.\par\end{minipage}
\end{table}

% Study 1: the differently worded effort item set aside.
% Written by tools/make_tables.py from the pipeline table(s) below (paths relative to
% Analysis/R Analysis Pipeline, each with the SHA-256 it was read with); do not edit by hand.
%   Study 1/tables/study1_item_sensitivity.csv  434b75e3bb2ffaac902aa9c073a2424283507dcb489a9e201fb97be7639f44e5
\begin{table}[htbp]
\centering
\small
\caption[Study 1: the differently worded effort item set aside]{\textbf{Study~1: the differently worded effort item set aside.} One effort item, for the first image with a figure, labeled its top two answers differently from the other nine. The table gives H1 on effort and H2 recomputed without that item's ratings, and H2 without that image on either rating. The registration names this as a sensitivity check, settled after the September 2026 estimates had been read (registered secondary; no verdict).}\label{tab:supp_item_sensitivity}
\setlength{\tabcolsep}{4pt}
\begin{tabular}{@{}>{\raggedright\arraybackslash}p{10.2cm}ll@{}}
\toprule
Ratings left out & Estimate (SE) & 95\% CI\\
\midrule
\multicolumn{3}{@{}l}{\textit{H1: the penalty on effort}}\\
None & 0.92 (0.15) & 0.62 to 1.22\\
The effort rating of the first image with a figure & 0.93 (0.15) & 0.63 to 1.22\\
\addlinespace[3pt]
\multicolumn{3}{@{}l}{\textit{H2: the penalty on effort minus the penalty on creativity}}\\
None & 0.47 (0.10) & 0.28 to 0.67\\
The effort rating of the first image with a figure & 0.48 (0.10) & 0.29 to 0.67\\
The first image with a figure, on both ratings & 0.47 (0.10) & 0.27 to 0.66\\
\bottomrule
\end{tabular}
\par\smallskip
\begin{minipage}{\textwidth}\footnotesize\raggedright Estimates are in points on the 1-to-5 scale, averaged with equal weight over images with and without a human figure; the penalty is the Human-label mean minus the AI-label mean. SE, standard error; CI, confidence interval. On that item the top two answers read `A lot of effort involved' and `Significant effort involved'; on the other nine, `Significant effort involved' and `A great deal of effort involved'. Every item is scored 1 to 5 by the position of the answer chosen.\par\end{minipage}
\end{table}


Supplementary Table~\ref{tab:supp_bootstrap_d} gives the registered effect sizes. Cohen's $d$ on participants' mean ratings was \SOneDCreat{} on creativity (95\% CI \SOneDCreatLo{} to \SOneDCreatHi) and \SOneDEff{} on effort (\SOneDEffLo{} to \SOneDEffHi). Averaging each person's ten ratings first removes much of the variation, so these values cannot be read against the usual small, medium and large benchmarks. The model-based values, \SOneDModelCreat{} and \SOneDModelEff, put the penalty on the scale of a single rating and are the ones to compare with other studies.

% Study 1: effect sizes.
% Written by tools/make_tables.py from the pipeline table(s) below (paths relative to
% Analysis/R Analysis Pipeline, each with the SHA-256 it was read with); do not edit by hand.
%   Study 1/tables/study1_bootstrap_d.csv  a223905d69ebbaaa35cb1f512a1a1cdefca4c5678b8057acd613e38876f048ca
\begin{table}[htbp]
\centering
\small
\caption[Study 1: effect sizes]{\textbf{Study~1: effect sizes of the label's penalty.} The registered effect sizes (registered secondary; no verdict). Averaging each person's ten ratings first removes much of the variation, so $d$ on participants' mean ratings is larger than the model-based $d$ by construction and cannot be read against the usual small, medium and large benchmarks. The model-based $d$ is the one to compare with other studies.}\label{tab:supp_bootstrap_d}
\setlength{\tabcolsep}{4pt}
\begin{tabular}{@{}lccc@{}}
\toprule
Rating & $d$ on participants' means & 95\% CI & Model-based $d$\\
\midrule
Creativity & 0.66 & 0.30 to 1.06 & 0.41\\
Effort & 1.15 & 0.77 to 1.64 & 0.80\\
\bottomrule
\end{tabular}
\par\smallskip
\begin{minipage}{\textwidth}\footnotesize\raggedright Cohen's $d$ here is the Human-label mean minus the AI-label mean, divided by the pooled standard deviation of participants' mean ratings (56 Human-label and 55 AI-label participants). Its 95\% confidence interval (CI) comes from 10{,}000 resamples of participants drawn within each label group. The model-based $d$ divides the H1 contrast by the square root of the primary model's summed variances (between participants, between images, and the rest), which puts it on the scale of a single rating.\par\end{minipage}
\end{table}


\paragraph*{The written explanations.}
Section~\ref{sec:suppA-text} describes the eight word lists that measured the explanations and how a script applied them without seeing the label or the rating. Every explanation was written after its rating, so none of these results shows why anyone rated as they did. Supplementary Table~\ref{tab:supp_table_2} gives how often each list flagged the Human group's explanations, and Supplementary Table~\ref{tab:supp_table_3} how often the effort explanations of each group mentioned prompting or machines; Supplementary Note~7 reads both. Supplementary Table~\ref{tab:suppA-denotation} gives the five registered word-meaning tests, an exploratory set whose corrected results are descriptive readings, not verdicts. Supplementary Table~\ref{tab:supp_denotation_slope_sensitivity} refits them with each test's difference free to vary from person to person, a check that the registration names and that changes no reading: every estimate kept its sign, and concept against prior exposure stayed short of significance.

% Study 1: what the Human-label explanations mentioned.
% Written by tools/make_tables.py from the pipeline table(s) below (paths relative to
% Analysis/R Analysis Pipeline, each with the SHA-256 it was read with); do not edit by hand.
%   Study 1/tables/study1_text_effort_prevalence.csv  e312979f6f455d06c2e566ffaa664e76acc741d5ffef4debcf986730f52832ac
%   Study 1/tables/study1_text_creativity_prevalence.csv  b4017bff3c84c961dd6dfc4ffccf4beb8b439a0b6e3e770c38756ab56ef81726
\begin{table}[htbp]
\centering
\small
\caption[Study 1: what the Human-label explanations mentioned]{\textbf{Study~1: what the Human-label group's explanations mentioned.} The share of that group's explanations of each rating that each word list flagged (descriptive). One explanation can mention several things, so the shares need not sum to 100\%.}\label{tab:supp_table_2}
\setlength{\tabcolsep}{4pt}
\begin{tabular}{@{}>{\raggedright\arraybackslash}p{8.6cm}rl@{}}
\toprule
Word list & Share & 95\% CI\\
\midrule
\multicolumn{3}{@{}l}{\textit{Effort explanations}}\\
Visible detail in the picture & 62.5\% & 56.4\% to 68.8\%\\
Time and industriousness & 10.0\% & 6.6\% to 13.8\%\\
Skill and training & 5.2\% & 2.1\% to 9.1\%\\
None of these three & 32.0\% & 26.8\% to 37.3\%\\
\addlinespace[3pt]
\multicolumn{3}{@{}l}{\textit{Creativity explanations}}\\
The idea behind the work & 26.8\% & 22.0\% to 31.8\%\\
Says the work looks familiar & 17.0\% & 13.0\% to 21.6\%\\
Says the work does not look familiar & 7.5\% & 5.0\% to 10.2\%\\
Raises familiarity either way & 24.5\% & 18.9\% to 30.4\%\\
Visible detail and rendering & 12.1\% & 9.1\% to 15.5\%\\
None of the idea, familiarity or detail & 47.1\% & 40.7\% to 53.9\%\\
\bottomrule
\end{tabular}
\par\smallskip
\begin{minipage}{\textwidth}\footnotesize\raggedright Each share is of the 560 explanations that the 56 Human-label participants wrote for that rating. CI, confidence interval, from 2{,}000 resamples of participants. A word list is a fixed list of words and phrases, applied by a script that saw neither the label nor the rating (Section~\ref{sec:suppA-text}).\par\end{minipage}
\end{table}

% Study 1: effort explanations that mention prompting or machines.
% Written by tools/make_tables.py from the pipeline table(s) below (paths relative to
% Analysis/R Analysis Pipeline, each with the SHA-256 it was read with); do not edit by hand.
%   Study 1/tables/study1_text_prompt_by_arm.csv  90cb709fae69a3f0623a99565fc0f97ed61b6d082321c6429757617775a2d858
\begin{table}[htbp]
\centering
\small
\caption[Study 1: effort explanations that mention prompting or machines]{\textbf{Study~1: effort explanations that mention prompting or machines, by label} (descriptive). Under the AI label such talk restates what participants were told. Under the Human label a mention is not a rejection of the account, since it can compare, deny or wonder.}\label{tab:supp_table_3}
\setlength{\tabcolsep}{4pt}
\begin{tabular}{@{}lrl@{}}
\toprule
Label & Share of effort explanations & 95\% CI\\
\midrule
Human label & 15.4\% & 10.4\% to 20.2\%\\
AI label & 50.5\% & 40.7\% to 60.3\%\\
\bottomrule
\end{tabular}
\par\smallskip
\begin{minipage}{\textwidth}\footnotesize\raggedright The share of each group's effort explanations that the prompting word list flagged: any mention of a machine, a prompt or a generated image. CI, confidence interval, from 2{,}000 resamples of participants. The test of prompting talk on effort against creativity, under the AI label, is in Supplementary Table~\ref{tab:suppA-denotation}.\par\end{minipage}
\end{table}

% Study 1's five registered word-meaning tests (an exploratory set, Benjamini-Hochberg corrected),
% every number a macro from generated/numbers.tex (source: Study 1/tables/study1_text_denotation_tests.csv).
% See supplement/tables/README_suppA.md.
\begin{table}[tbp]
\centering
\small
\caption{\textbf{Study~1's five word-meaning tests.} Each test asks whether one word list flagged explanations more often than another, in the explanations named, or, in the last row, whether one list flagged effort explanations more often than creativity explanations. Each is a mixed logistic regression, a model for a yes-or-no outcome (whether a list flagged an explanation) with a baseline for each participant and each image, so that one person's twenty explanations do not count as twenty independent observations. The estimate is the difference in the log-odds that the first list flags an explanation rather than the second; a positive value is the predicted direction, and all five are positive. The registration names the five as one exploratory set corrected by the Benjamini--Hochberg method, which limits the expected share of false positives among results called significant; the corrected results are descriptive readings, not verdicts. The detail-against-skill fit converged only when restarted from where its first attempt stopped; the prompting fit was singular, with no variation between images. Supplementary Note~7 gives the shares of explanations behind each test.}
\label{tab:suppA-denotation}
\begin{tabular}{@{}>{\raggedright\arraybackslash}p{4.4cm}>{\raggedright\arraybackslash}p{3.6cm}llll@{}}
\toprule
Explanations & Comparison & Estimate & SE & Uncorrected & Corrected\\
\midrule
Effort, Human label & Visible detail against time & \SOneTextDetailVsTimeEst & \SOneTextDetailVsTimeSE & $p$~\SOneTextDetailVsTimeP & $p$~\SOneTextDetailVsTimePBH\\
Effort, Human label & Visible detail against skill & \SOneTextDetailVsSkillEst & \SOneTextDetailVsSkillSE & $p$~\SOneTextDetailVsSkillP & $p$~\SOneTextDetailVsSkillPBH\\
Creativity, Human label & Concept against visible detail (execution) & \SOneTextConceptVsExecutionEst & \SOneTextConceptVsExecutionSE & $p$~\SOneTextConceptVsExecutionP & $p$~\SOneTextConceptVsExecutionPBH\\
Creativity, Human label & Concept against prior exposure (novelty) & \SOneTextConceptVsNoveltyEst & \SOneTextConceptVsNoveltySE & $p$~\SOneTextConceptVsNoveltyP & $p$~\SOneTextConceptVsNoveltyPBH\\
AI label, effort against creativity & Prompting or machines & \SOneTextPromptEffortVsCreatEst & \SOneTextPromptEffortVsCreatSE & $p$~\SOneTextPromptEffortVsCreatP & $p$~\SOneTextPromptEffortVsCreatPBH\\
\bottomrule
\end{tabular}
\end{table}

% Study 1: the word-meaning tests with a random slope.
% Written by tools/make_tables.py from the pipeline table(s) below (paths relative to
% Analysis/R Analysis Pipeline, each with the SHA-256 it was read with); do not edit by hand.
%   Study 1/tables/study1_text_denotation_slope_sensitivity.csv  7b92f279a83a0c66990ce6bf3c8a500da9ce97b00e1cab45c8c7fb412802875a
\begin{table}[htbp]
\centering
\footnotesize
\caption[Study 1: the word-meaning tests with a random slope]{\textbf{Study~1: the five word-meaning tests with each person's difference free to vary.} The registered model gives each participant a baseline of their own but one difference between the two lists that all participants share; the refit lets that difference vary from person to person. The registration names this refit as a sensitivity check that cannot change a reading (exploratory; its $p$-values are uncorrected). Every estimate kept its sign.}\label{tab:supp_denotation_slope_sensitivity}
\setlength{\tabcolsep}{4pt}
\begin{tabular}{@{}>{\raggedright\arraybackslash}p{3.9cm}>{\raggedright\arraybackslash}p{5.0cm}llrl@{}}
\toprule
Explanations & Comparison & Registered & Refit & $z$ & $p$\\
\midrule
Effort, Human label & Visible detail against time & 3.08 (0.19) & 3.57 (0.37) & 9.60 & $p$~\ensuremath{<}~.001\\
Effort, Human label & Visible detail against skill & 3.87 (0.24) & 5.66 (0.95) & 5.98 & $p$~\ensuremath{<}~.001\\
Creativity, Human label & Concept against visible detail (execution) & 1.03 (0.16) & 1.03 (0.22) & 4.63 & $p$~\ensuremath{<}~.001\\
Creativity, Human label & Concept against prior exposure (novelty) & 0.13 (0.14) & 0.28 (0.21) & 1.30 & $p$~\ensuremath{=}~.195\\
AI label, effort against creativity & Prompting or machines & 1.78 (0.18) & 2.11 (0.46) & 4.58 & $p$~\ensuremath{<}~.001\\
\bottomrule
\end{tabular}
\par\smallskip
\begin{minipage}{\textwidth}\footnotesize\raggedright Each cell gives an estimate with its standard error (SE) in parentheses, from the registered model and from the refit. Estimates are differences in log-odds, the logarithm of the odds that the first list flags an explanation rather than the second; a positive value is the predicted direction. $z$ and $p$ test the refit's estimate against zero. Every refit converged and none was singular (a variance or correlation estimated at the edge of its range). Supplementary Table~\ref{tab:suppA-denotation} gives the registered tests with their corrected $p$-values.\par\end{minipage}
\end{table}


The registration also asks whether a participant rated lower, on creativity, the particular images that they called familiar (registered secondary, description only). Within a person, an image whose explanation called it familiar received a creativity rating that differed from that person's other images by \SOneFamiliarCreatHumanEst{} points under the Human label (95\% CI \SOneFamiliarCreatHumanLo{} to \SOneFamiliarCreatHumanHi) and by \SOneFamiliarCreatAIEst{} under the AI label (\SOneFamiliarCreatAILo{} to \SOneFamiliarCreatAIHi); a negative value means a lower rating. The registration names effort as the point of comparison; there the values were \SOneFamiliarEffHumanEst{} (\SOneFamiliarEffHumanLo{} to \SOneFamiliarEffHumanHi) and \SOneFamiliarEffAIEst{} (\SOneFamiliarEffAILo{} to \SOneFamiliarEffAIHi). The explanation followed the rating, so these associations do not show that calling an image familiar lowered its rating. The analysis table \nolinkurl{study1_text_novelty_within_between.csv} gives every term.

The two checks run in every study (Section~\ref{sec:suppA-text}) are registered secondary. In the first, the words of explanations from participants left out of the fit predicted the rating, with correlations of \SOneTextRecRatingCreat{} on creativity (95\% CI \SOneTextRecRatingCreatLo{} to \SOneTextRecRatingCreatHi) and \SOneTextRecRatingEff{} on effort (\SOneTextRecRatingEffLo{} to \SOneTextRecRatingEffHi). Both lie above the 95th percentile of the correlations obtained after each person's ratings were shuffled among that person's own explanations (\SOneTextRecRatingCreatNull{} and \SOneTextRecRatingEffNull). They also told the label groups apart. The AUC, the chance that a randomly chosen AI-label explanation scores above a randomly chosen Human-label one, so that 0.5 is chance, was \SOneTextRecArmCreat{} on creativity (\SOneTextRecArmCreatLo{} to \SOneTextRecArmCreatHi) and \SOneTextRecArmEff{} on effort (\SOneTextRecArmEffLo{} to \SOneTextRecArmEffHi). Added to the rating, the words raised the AUC by \SOneTextRecIncrCreat{} (\SOneTextRecIncrCreatLo{} to \SOneTextRecIncrCreatHi) and \SOneTextRecIncrEff{} (\SOneTextRecIncrEffLo{} to \SOneTextRecIncrEffHi), so by the registered reading the words add to the rating in telling the groups apart. Recovery of the rating was already partly known from earlier exploratory work (Section~\ref{sec:suppA-registrations}). In the second check, AI-label participants reused their own words across images more than Human-label participants on effort (difference \SOneTextOverlapEffDiff, 95\% CI \SOneTextOverlapEffLo{} to \SOneTextOverlapEffHi; Holm $p$~\SOneTextOverlapEffPHolm), as predicted, but not detectably on creativity (\SOneTextOverlapCreatDiff, \SOneTextOverlapCreatLo{} to \SOneTextOverlapCreatHi; Holm $p$~\SOneTextOverlapCreatPHolm). The registered reading is `supported' on effort and `inconclusive' on creativity, which is not evidence that the groups reuse their words equally. Supplementary Tables~\ref{tab:supp_text_recoverability} and~\ref{tab:supp_text_selfsimilarity} give every row.

% Study 1: the rating and the label recovered from the words.
% Written by tools/make_tables.py from the pipeline table(s) below (paths relative to
% Analysis/R Analysis Pipeline, each with the SHA-256 it was read with); do not edit by hand.
%   Study 1/tables/study1_text_recoverability.csv  23af11956d87fa145255dddef8f9a3ded2c67aedc1e6c04a61a6879708e0babe
\begin{table}[htbp]
\centering
\small
\caption[Study 1: the rating and the label recovered from the words]{\textbf{Study~1: can the rating and the label be recovered from the words?} The first of the two checks run in every study (registered secondary). Each model was fitted to some participants' explanations and scored on the explanations of participants left out of the fit. By the registered reading, both the rating and the label are recovered on both ratings, and the words add to the rating in telling the two label groups apart.}\label{tab:supp_text_recoverability}
\setlength{\tabcolsep}{4pt}
\begin{tabular}{@{}>{\raggedright\arraybackslash}p{6.4cm}llll@{}}
\toprule
What is predicted & Statistic & Estimate & 95\% CI & Shuffled baseline\\
\midrule
\multicolumn{5}{@{}l}{\textit{Creativity explanations}}\\
The rating, from the words & $r$ & 0.49 & 0.44 to 0.55 & 0.16\\
The label, from the words & AUC & 0.59 & 0.53 to 0.64 & 0.51\\
The label, from the rating & AUC & 0.60 & 0.53 to 0.66 & 0.51\\
The label, from the rating and the words & AUC & 0.63 & 0.57 to 0.69 & --\\
What the words add: rating and words, minus rating alone & AUC difference & 0.036 & 0.003 to 0.067 & --\\
\addlinespace[3pt]
\multicolumn{5}{@{}l}{\textit{Effort explanations}}\\
The rating, from the words & $r$ & 0.63 & 0.57 to 0.68 & 0.39\\
The label, from the words & AUC & 0.72 & 0.66 to 0.78 & 0.50\\
The label, from the rating & AUC & 0.70 & 0.63 to 0.77 & 0.53\\
The label, from the rating and the words & AUC & 0.75 & 0.68 to 0.81 & --\\
What the words add: rating and words, minus rating alone & AUC difference & 0.051 & 0.020 to 0.083 & --\\
\bottomrule
\end{tabular}
\par\smallskip
\begin{minipage}{\textwidth}\footnotesize\raggedright $r$, the correlation between predicted and observed ratings (0, no relation; 1, a perfect one). AUC, the chance that a randomly chosen AI-label explanation is scored as more AI-like than a randomly chosen Human-label one (0.5 is chance; 1, perfect separation). The models are ridge regressions, which shrink their many word weights toward zero, on counts of each explanation's content words. Participants were split into 10 groups, and each group in turn was left out of the fit and scored, so that no participant's explanations were both fitted and scored. CI, confidence interval, from 2{,}000 resamples of participants. The shuffled baseline is the 95th percentile of the statistic over 200 reruns with the outcome shuffled: the rating among each participant's own explanations, the label across participants. An estimate above it is read as signal. Every explanation was written after its rating, so recovering the rating is not evidence that the words produced it.\par\end{minipage}
\end{table}

% Study 1: reuse of one's own words, by label.
% Written by tools/make_tables.py from the pipeline table(s) below (paths relative to
% Analysis/R Analysis Pipeline, each with the SHA-256 it was read with); do not edit by hand.
%   Study 1/tables/study1_text_selfsimilarity_by_arm.csv  1503b1471e4774488e277de4b26e8415543040f5c3d3897169696890288e0672
\begin{table}[htbp]
\centering
\small
\caption[Study 1: reuse of one's own words, by label]{\textbf{Study~1: how much each person reused their own words, by label.} The second of the two checks run in every study (registered secondary). It predicted that participants shown the AI label would reuse their own words more. By the registered reading the prediction is supported on effort and inconclusive on creativity, which is not evidence that the two groups reuse their words equally.}\label{tab:supp_text_selfsimilarity}
\setlength{\tabcolsep}{4pt}
\begin{tabular}{@{}>{\raggedright\arraybackslash}p{7.6cm}ll@{}}
\toprule
Measure & Human label (95\% CI) & AI label (95\% CI)\\
\midrule
\multicolumn{3}{@{}l}{\textit{Creativity explanations}}\\
Overlap among a person's own explanations & 0.069 (0.055 to 0.083) & 0.094 (0.070 to 0.124)\\
Overlap with others' explanations of the same image & 0.027 (0.024 to 0.031) & 0.029 (0.026 to 0.033)\\
Own overlap minus others' overlap & 0.041 (0.029 to 0.054) & 0.065 (0.041 to 0.094)\\
Content words per explanation & 8.03 (6.80 to 9.40) & 9.17 (7.88 to 10.55)\\
\multicolumn{3}{@{}p{\textwidth}@{}}{\raggedright Registered comparison (AI minus Human in own minus others' overlap, at equal length): 0.025 (95\% CI \ensuremath{-}0.003 to 0.059); $p$~\ensuremath{=}~.092; Holm $p$~\ensuremath{=}~.092; inconclusive.}\\
\addlinespace[3pt]
\multicolumn{3}{@{}l}{\textit{Effort explanations}}\\
Overlap among a person's own explanations & 0.099 (0.082 to 0.118) & 0.157 (0.121 to 0.201)\\
Overlap with others' explanations of the same image & 0.036 (0.032 to 0.041) & 0.036 (0.030 to 0.042)\\
Own overlap minus others' overlap & 0.063 (0.046 to 0.081) & 0.121 (0.088 to 0.163)\\
Content words per explanation & 7.99 (6.63 to 9.53) & 7.89 (6.76 to 9.16)\\
\multicolumn{3}{@{}p{\textwidth}@{}}{\raggedright Registered comparison (AI minus Human in own minus others' overlap, at equal length): 0.057 (95\% CI 0.021 to 0.104); $p$~\ensuremath{<}~.001; Holm $p$~\ensuremath{=}~.002; supported.}\\
\bottomrule
\end{tabular}
\par\smallskip
\begin{minipage}{\textwidth}\footnotesize\raggedright Overlap is the share of content words that two explanations have in common: the words in both divided by the words in either (0, none in common; 1, the same words). A person's own overlap is averaged over all pairs of that person's explanations of one rating; others' overlap pairs each of those explanations with other participants' explanations of the same image and rating. The registered comparison is a regression of own minus others' overlap on the label, holding the mean number of content words constant, because overlap falls as texts get longer. CI, confidence interval, from 2{,}000 resamples of participants within each label group, with others' overlap rebuilt each time. Holm's correction runs over the two ratings. Supported means that the corrected test rejected no difference, in the predicted direction; inconclusive, that it did not.\par\end{minipage}
\end{table}


The registration names the general description of the explanations as exploratory, and none of it carries a verdict. The tone and concreteness of the words, scored from published norms, did not differ between the label groups on either rating after a Benjamini--Hochberg correction, which limits the expected share of false positives among results called significant (analysis table \nolinkurl{study1_text_affect_by_arm.csv}). A topic model, a method that finds groups of words that tend to occur together without being told the label or the rating, chose the number of topics for each rating by a stability rule (Supplementary Table~\ref{tab:supp_text_topics_stability}). After correction, AI-label explanations gave more weight on both ratings to a topic led by `ai', with `created' and `prompt' among its leading words on creativity and `generated' and `generator' on effort. On effort they gave less weight to a topic led by `details', `small' and `many' and to one led by `looks', `generated', `good' and `real' (analysis table \nolinkurl{study1_text_topics_by_arm.csv}). The topics agreed only weakly with groups formed from sentence embeddings, lists of numbers that a pretrained model assigns to each explanation (analysis table \nolinkurl{text_topics_embedding_agreement.csv}). A lasso, a regression that sets most word weights to exactly zero, kept few words for telling the label groups apart once the rating was known, led by `due', `prompt' and `ai' (analysis table \nolinkurl{study1_text_prediction_weights.csv}).

% Study 1: how the number of topics was chosen.
% Written by tools/make_tables.py from the pipeline table(s) below (paths relative to
% Analysis/R Analysis Pipeline, each with the SHA-256 it was read with); do not edit by hand.
%   Study 1/tables/text_topics_stability.csv  fa312640420b5534bc1eceb367d6f094fbd1949ca6765c465870583d9c140475
\begin{table}[htbp]
\centering
\small
\caption[Study 1: how the number of topics was chosen]{\textbf{Study~1: how the number of topics was chosen.} For each number of topics from 3 to 10, how stable the topic model of each rating's explanations was (exploratory). The most stable number was used: 6 for both ratings (in bold).}\label{tab:supp_text_topics_stability}
\setlength{\tabcolsep}{4pt}
\begin{tabular}{@{}lll@{}}
\toprule
Number of topics & Creativity: stability (SD) & Effort: stability (SD)\\
\midrule
3 & 0.83 (0.13) & 0.86 (0.15)\\
4 & 0.89 (0.08) & 0.87 (0.14)\\
5 & 0.88 (0.07) & 0.94 (0.08)\\
6 & \textbf{0.92 (0.05)} & \textbf{0.98 (0.01)}\\
7 & 0.84 (0.07) & 0.95 (0.05)\\
8 & 0.82 (0.08) & 0.91 (0.06)\\
9 & 0.84 (0.06) & 0.87 (0.06)\\
10 & 0.87 (0.07) & 0.86 (0.06)\\
\bottomrule
\end{tabular}
\par\smallskip
\begin{minipage}{\textwidth}\footnotesize\raggedright A topic model (here non-negative matrix factorization) finds groups of words that tend to occur together. At each number of topics it was refitted ten times, each time on a random 80\% of the participants. The topics of two refits were matched one to one, and their agreement is the mean cosine similarity of the matched topics' word weights (1 when two topics weight the same words in the same proportions, 0 when they share no word). Stability is that agreement averaged over every pair of refits; SD, its standard deviation. A tie at three decimals goes to the smaller number.\par\end{minipage}
\end{table}

\endgroup
